\begin{theorem}\nf{(Признак Харди)}\\
	Пусть $f\in\mc R[a,b],\ a<b<+\infty$ и периодична с периодом $\omega$, $g$ монотонна на $[a,+\infty)$ и бесконечно мала при $x\ra+\infty$. Тогда:
	\begin{enumerate}
		\item $\ds\int_a^{a+\omega}f(x)dx=0$, то $\ds\int_a^{+\infty}f(x)g(x)dx$ сходится.
		\item $\ds\int_a^{a+\omega}f(x)dx\neq0$, то $\ds\int_a^{+\infty}f(x)g(x)dx$ и $\ds\int_a^{+\infty}g(x)dx$ сходятся.
	\end{enumerate}
\end{theorem}
\begin{proof}
	Вначале докажем первый пункт.
	\[(\fa k\in\N)\ \int_a^{a+k\cdot\omega}f(x)dx=\sum_{i=1}^k\int_{a+(i-1)\cdot\omega}^{a+i\cdot\omega}f(x)dx=\sum_{i=1}^k\int_a^{a+\omega}f(x)dx=0\]
	Заметим: $(\fa A>a)(\ex k\in\N)\ a+(k-1)\cdot\omega<A\le a+k\cdot\omega$\\
	Тогда:
	\begin{multline*}
		\l|\int_a^Af(x)dx\r|=
		\l|\int_a^{a+(k-1)\cdot\omega}f(x)dx+\int_{a+(k-1)\cdot\omega}^Af(x)dx\r|=
		\Bigg|\int_{a+(k-1)\cdot\omega}^Af(x)dx\Bigg|=\\=
		\l|\int_a^{A-(k-1)\cdot\omega}f(x)dx\r|\le
		\Bigg|\int_a^{a+\omega}f(t)dt\Bigg|=:M
	\end{multline*}
	Получили требуемое.\\
	Теперь второй пункт.\\
	$\ds K:=\int_a^{a+\omega}f(x)dx\neq0,\ f_1(x):=f(x)-\frac K\omega$\\
	\[\int_a^{a+\omega}f_1(x)dx=\int_a^{a+\omega}f(x)dx-\frac K\omega\cdot\omega=0\]
	Теперь можем воспользоваться первым пунктом для $f_1$: $\ds\int_a^{+\infty}f_1(x)g(x)dx$ сходится.\\
	При этом знаем: $\ds\int_a^{+\infty}f_1(x)g(x)dx=\int_a^{+\infty}f(x)g(x)dx-\frac K\omega\int_a^{+\infty}g(x)dx$
	Получили требуемое.
\end{proof}
\begin{examples}
	Рассмотрим $\ds\int_1^\p e^{\cos x}\sin(\sin x)\cdot\frac{dx}x\text{ и }
	\int_1^\p e^{\sin x}\cdot\sin(\sin x)\cdot\frac{dx}x$
	\[\int_0^{2\pi}e^{\cos x}\sin(\sin x)dx=\l(\int_0^\pi+\int_\pi^{2\pi}\r)e^{\cos x}\sin(\sin x)dx=0\]
	\[\int_0^{2\pi}e^{\sin x}\sin(\sin x)dx=\int_0^\pi e^{\sin x}\sin(\sin x)dx-\int_0^\pi e^{-\sin x}\sin(\sin x)dx>0\]
\end{examples}
\begin{example}
	Теперь покажем, что заменой на эквивалентные функции можно пользоваться только в случае, если функции знакопостоянны.\\
	Рассмотрим $\ds\int_2^\p \frac{\sin x}{x^\al+\sin x}dx$
	\[\frac{\sin x}{x^\al+\sin x}=\frac{\sin x}{x^\al}-\frac{\sin^2x}{x^\al(x^\al+\sin x)}\]
	$\ds\frac{\sin x}{x^\al}$ уже был ранее исследован, при 
	\[\frac{\sin^2x}{x^\al(x^\al+1)}\le\frac{\sin^2x}{x^\al(x^\al+\sin x)}\le\frac{\sin^2x}{x^\al(x^\al-1)}\le\frac1{x^\al(x^\al-1)}\sim\frac1{x^{2\al}}\]
	$2\al>1,\ \al>\frac12$. При $\al\le\frac12$ эта штука расходится.\\
	Заметим, что если бы мы воспользовались той штукой, которую я не помню как зовут, то:
	\[\frac{\sin x}{x^\al+\sin x}=\frac{\sin x}{x^\al\cdot(1+\frac{\sin x}{x^\al})}\sim\frac{\sin x}{x^\al}\]
	А уже последняя полученная штука чуть выше сходится при $\al>0$, неправильный ответ.
\end{example}
\section{Теория рядов}
\subsection{Числовые ряды}
\begin{definition}
	Скажем формально, далее поймём что это значит.\\
	\textit{Числовым рядом} называется сумма бесконечного числа чисел.\\ Обозначается как $\ds\sum_{i=1}^\infty a_n$. \textit{Частичной суммой ряда} называется $\ds S_n:=\sum_{k=1}^na_k$.\\
	Ряд сходится, если $\ds\{S_n\}_{n=1}^\infty$ имеет конечный предел.
\end{definition}
\begin{lemma}\nf{(Элементарные свойства рядов)}\\
	\begin{enumerate}
		\item Если ряды $\ds\sum_{n=1}^\i a_n$ и $\ds\sum_{n=1}^\i b_n$ сходятся, то ряд $\ds(\fa\al,\beta\in\R)\ \sum_{n=1}^\i(\al a_n+\beta b_n)$ сходится, причём $\ds\sum_{n=1}^\i(\al a_n+\beta b_n)=\al\sum_{n=1}^\i a_n+\beta\sum_{n=1}^\i b_n$
		\item $(\fa n_0\in\N)$ ряды $\ds\sum_{n=1}^\i$ и $\ds\sum_{n=n_0}^\i$ сходятся или расходятся одновременно.
	\end{enumerate}
\end{lemma}
\begin{proof}
	Доказательство предлагаем проделать читателю в качестве упражнения.
\end{proof}
\begin{theorem}\nf{(Необходимое условие сходимости ряда)}\label{12_th1}\\
	Если ряд $\ds\sum_{n=1}^\i a_n$ сходится, то $\lim_{n\ra\p}a_n=0$
\end{theorem}
\begin{proof}
	По условию: $\ds\ex\lim_{n\ra\p}S_n=S$\\
	Тогда: $a_n=S_n-S_{n-1}$. При $n\ra\p$: $a_n\ra S-S=0$.\\
	Получили требуемое.
\end{proof}
\begin{example}
	Рассмотрим как пример геометрическую прогрессию $\ds\sum_{n=1}^\i q_n$.
	\begin{enumerate}
		\item $|q|\ge1\Ra\nexists\lim_{n\ra\p}q_n$, расходится.
		\item $|q|<1$. $S_n=q\cdot\frac{1-q^n}{1-q}\ra\frac q{1-q}=\sum_{n=1}^\i q^n$ при $n\ra\p$.
	\end{enumerate}
\end{example}
\begin{theorem}\nf{(Критерий сходимости рядов с неотрицательными членами)}\label{12_th2}\\
	Если $a_n\ge0\ \fa n\in\N$, то $\ds\sum_{n=1}^\i a_n$ сходится тогда и только тогда, когда последовательность его частичных сумм ограничена.
\end{theorem}
\begin{proof}
	$\ds a_n\ge0\Ra S_{n-1}\le S_n\Ra\{S_n\}_{n=1}^\i$ возрастает $\ds\Ra\lim_{n\ra\p}S_n=\sup\{S_n\}$\\
	Получили требуемое.
\end{proof}
\begin{theorem}\label{12_th3}\nf{(Признак сходимости)}\\
	Если $0\le a_n\le b_n\ \fa n\in\N$, то:
	\begin{enumerate}
		\item $\ds\ss b_n$ сходится, тогда $\ds\ss a_n$ сходится.
		\item $\ds\ss a_n$ расходится, тогда $\ds\ss b_n$ расходится.
	\end{enumerate}
\end{theorem}
\begin{proof}
	$A_n,B_n$ - соответсвующие частичные суммы.\\
	$\ds A_n=\sum_{k=1}^n a_k\le B_n=\sum_{k=1}^n b_k$\\
	Получить требуемое можно из выражения выше.
\end{proof}
\begin{note}
	Внимательный читатель мог заметить некоторое сходство между теоремой (\ref{12_th3}) и примерно такой же для несобственных интегралов. Между этими штуками есть связь, эта секретная тайна загадка будет раскрыта в следующей теореме.
\end{note}
\begin{theorem}\label{12_th4}\nf{(Интегральный признак Коши - Маклорена)}\\
	Если $f(x)\ge0$ и $f$ убывает на $[1,\p]$, то $\ds\ss f(n)$ и $\ds\int_1^pf(x)dx$ сходятся и расходятся одновременно.
\end{theorem}
\begin{proof}
	\[\ds\fa x\in[k,k+1]\ f(k)\ge f(x)\ge f(k+1),\ f(k)\ge\int_1^{n+1}f(x)dx\ge f(k+1)\]
	\[\sum_{k=1}^nf(k)\ge\int_1^{n+1}f(x)dx\ge\sum_{k=1}^nf(k+1)=\sum_{j=2}^{n+1}f(j)\]
	Получили требуемое.
\end{proof}
\begin{example}
	Пусть $\ds f(x)=\frac1{x^\al}\ge0,\ \ss\frac1{x^\al}$\\
	$\int_a^\p\frac{dx}{x^\al}$ сходится при $\al>1$, значит и $\ss\frac1{x^\al}$ тоже сходится.
\end{example}
\begin{theorem}\label{12_th5}\nf{(Признак Даламбера)}\\
	Пусть $a_k>0\ \fa k\in\N$. Тогда:
	\begin{enumerate}
		\item $\ds\v\lim_{n\ra\i}\frac{a_{n+1}}{a_n}=p<1$, тогда ряд $\ds\ss a_n$ сходится.
		\item $\ds\underline\lim_{n\ra\i}\frac{a_{n+1}}{a_n}=q>1$, тогда ряд $\ds\ss a_n$ расходится.
	\end{enumerate}
\end{theorem}
\begin{proof}
	По определению верхнего предела:
	\[(\fa\eps>0)(\ex N\in\N)(\fa n\ge N)\ \frac{a_{n+1}}{a_n}<p+\eps=\t p<1\]
	\[(\fa k\in N)\ a_{N+k}=\frac{a_{N+k}}{a_{N+k-1}}\cdot\frac{a_{N+k-1}}{a_{N+k-2}}\cdot\dots\cdot\frac{a_{N+1}}{a_N}\cdot a_N<\t p^k\cdot a_N\]
	$\ds\ss[N+1]a_n$ сходится по (\ref{12_th3}).
	Получили требуемое.
\end{proof}
\begin{theorem}\label{12_th6}\nf{(Признак Коши)}\\
	Пусть $a_k>0\ \fa k\in\N$. Тогда:
	\begin{enumerate}
		\item $\ds\v{\lim_{n\ra\p}}\sqrt[n]{a_n}<1$, тогда $\ds\ss a_n$ сходится.
		\item $\ds\v{\lim_{n\ra\p}}\sqrt[n]{a_n}>1$, тогда $\ds\ss a_n$ расходится.
	\end{enumerate}
\end{theorem}
\begin{proof}
	Докажем первый пункт.
	\[p:=\v{\lim_{n\ra\p}}\sqrt[n]{a_n},\ (\fa\eps>0)(\ex N\in\N)(\fa n>N)\ \sqrt[n]{a_n}<p+\eps=\t p<1,\ a_n<(\t p)^n,\ \ss[N+1]\t p\text{ сходится}\]
	Теперь второй. $\ds\lim_{n\ra\p}\sqrt[n]{a_n}=p>1,\ \ex\{n_k\}_{k=1}^\i,\ \lim_{k\ra\p}\sqrt[n]{a_{n_k}}=p$
	\[(\fa\eps>0)(\ex K\in\N)(\fa k>K)\ \sqrt[n]{a_{n_k}}>p-\eps=q>1\Ra a_{n_k}>q^{n_k}\ra\p,\ k\ra\p,\ \lim_{k\ra\p}a_{n_k}\neq0\]
	Получили требуемое.
\end{proof}
\begin{note}
	Теоремы (\ref{12_th5}) и (\ref{12_th6}) иногда записывают в предельной форме. Отличие заключается в том, что предел становится не верхний, а обычный, и также добавляется условие: при $\lim=1$ ничего про сходимость с помощью этих теорем сказать нельзя.
\end{note}